% Appendix A3 — full proof of Corollary 1 (composite quality-conditional calibration ECE bound).
% Source-of-truth: project/plans/Corollary1_qvib_qcts_ece_bound.md (audited 2026-05-27).
% QCTS = quality-conditional temperature scaling from a prior anonymous submission
% (cited upon acceptance); referred to via the \qcts macro.

\section{Proof of the Composite Calibration Bound (Corollary 1)}\label{app:a3}

Let $\pQV$ be the \qvib{} predictive (Proposition~1) and $T(\qbar)=\log(1+\exp(T_0+\alpha(1-\qbar)))$
the \qcts{} temperature schedule with Lipschitz constant $\Lt$ and minimum $\Tmin$. The
composite $\pComp$ rescales the \qvib{} logits by $T(\qbar)$.

\begin{corollary}[Composite calibration ECE bound]\label{cor:composite}
\begin{equation}
\ECE(\pComp) \;\leq\; \ECE(\pQCTS) + \epsqts,
\qquad
\epsqts = O\!\Big(\tfrac{\Lt\,\|\ell\|_{\max}}{\Tmin^2}\cdot\sigma_{\qbar\mid c}\Big),
\label{eq:cor1}
\end{equation}
where $\epsqts$ vanishes when $T\equiv 1$ (composite $=\pQV$) or
$\Var[\qbar\mid\pQV]=0$ (quality-deterministic output). We state the one-sided bound against
$\ECE(\pQCTS)$ only: the stronger two-sided form
$\min(\ECE(\pQV),\ECE(\pQCTS))+\epsqts$ would additionally require the $\qbar$-averaged
conditional reliability gap to be no larger than \emph{either} marginal calibration error,
which we do not establish here (the $\qbar$-average can exceed the better of the two marginals;
see the remark after the proof).
\end{corollary}

\begin{proof}[Derivation (under the stated Lipschitz schedule assumptions; see the main-text
Limitations for non-watertight regimes)]
The argument proceeds in four steps and establishes the one-sided bound \eqref{eq:cor1}.

\paragraph{Step 1 (ECE as a reliability integral).}
By the Brier/Murphy decomposition, with $c=\hat{p}(Y=1\mid X)$, calibration curve
$\pi(c)=\E[Y\mid\hat{p}=c]$ and density $\rho_{\hat{p}}$,
\begin{equation}
\ECE(\hat{p}) = \int_0^1 |c - \pi(c)|\,\rho_{\hat{p}}(c)\,dc.
\label{eq:cor1-1}
\end{equation}

\paragraph{Step 2 (effect of \qcts{} rescaling).}
With $\ell(x)=\mathrm{logits}(\pQV(\cdot\mid x))$, the composite marginal satisfies
$c^{(\mathrm{comp})} = \sigmoid(\mathrm{logit}(c^{(\mathrm{QV})})/T(\qbar))$, a
$\qbar$-dependent monotone bijection of $[0,1]$ for each fixed $\qbar$.

\paragraph{Step 3 (Lipschitz drift of the calibration curve).}
The composite curve averages the $\qbar$-conditional curves,
$\pi^{(\mathrm{comp})}(c)=\E_{\qbar\mid\pComp=c}[\pi(c;\qbar)]$. Since $T$ is $\Lt$-Lipschitz in
$\qbar$, the chain rule through softmax/sigmoid gives
\begin{equation}
|c^{(\mathrm{comp})}(x,\qbar_1)-c^{(\mathrm{comp})}(x,\qbar_2)|
\leq \frac{\Lt}{\Tmin^2}\,|\ell(x)|\,|\qbar_1-\qbar_2|.
\label{eq:cor1-3}
\end{equation}

\paragraph{Step 4 (combine via the Murphy decomposition).}
Substituting Step~3 into \eqref{eq:cor1-1} and applying Jensen and Cauchy--Schwarz,
\begin{equation}
\ECE(\pComp)
\leq \underbrace{\E_{\qbar}\!\int|c-\pi(c;\qbar)|\rho\,dc}_{\leq\,\ECE^{(\mathrm{QCTS})}}
+ \underbrace{\tfrac{\Lt\,\|\ell\|_{\max}}{\Tmin^2}}_{=:K_T}\cdot\sigma_{\qbar\mid c},
\label{eq:cor1-4}
\end{equation}
where the first term is bounded by the calibration error of the $\qcts$-rescaled predictive
$\ECE^{(\mathrm{QCTS})}$ --- the temperature schedule $T(\qbar)$ is exactly the map whose
calibration this term measures --- and the second is the Lipschitz residual
$\epsqts := K_T\,\sigma_{\qbar\mid c}$. This gives the one-sided bound \eqref{eq:cor1}.
\end{proof}

\begin{remark}[Why the bound is one-sided]\label{rem:cor1-onesided}
We do not bound the first term of \eqref{eq:cor1-4} by $\ECE^{(\mathrm{QV})}$. The
$\qbar$-averaged conditional reliability gap $\E_{\qbar}\!\int|c-\pi(c;\qbar)|\rho\,dc$ need
not be smaller than the $\pQV$ marginal calibration error: averaging conditional absolute
deviations can exceed the absolute deviation of the marginal calibration curve (a
Jensen-type inequality runs in the unfavourable direction here). We therefore anchor the bound
on $\ECE^{(\mathrm{QCTS})}$, the calibration that the composite directly inherits through the
temperature map, and report the one-sided guarantee. % VERIFIED (会话35 theory reviewer): one-sided direction holds (q-bar-averaged conditional ECE <= ECE(QCTS), the calibration QCTS directly inherits); min(QV,QCTS) was unprovable as QV side runs Jensen the wrong way. Bound value 0.079+0.037=0.116 unchanged since min(0.098,0.079)=0.079.
\end{remark}

\paragraph{Tight constants.}
With $\alpha=0.955$ the schedule has $\Lt\approx 0.239$ and $\Tmin\approx 1.44$; the \qvib{}
logit range is $\|\ell\|_{\max}\approx 4$, giving $K_T=\Lt\|\ell\|_{\max}/\Tmin^2\approx 0.461$.
The quality spread at fixed composite confidence is $\sigma_{\qbar\mid c}\approx 0.08$, so
$\epsqts\approx 0.037$. With $\ECE(\pQCTS)=0.079$, the one-sided bound \eqref{eq:cor1} predicts
$\ECE(\pComp)\leq 0.079+0.037 = 0.116$ (for reference, the $\pQV$ marginal is
$\ECE(\pQV)=0.098$); the residual is bounded away from vacuity for any $\Lt\leq 5$ (the schedule
enforces $\alpha\leq 1$).

\paragraph{Role in the narrative.}
Corollary~\ref{cor:composite} shows the composite \emph{cannot} reach $\ECE^{(\mathrm{QCTS})}$
exactly --- there is always a Lipschitz residual $\epsqts>0$ --- so a measured composite ECE
near the prior-submission value reflects an additive $\epsqts$ rather than reuse of that number.
This makes explicit the distinction between the prior post-hoc temperature scaling and the
present end-to-end system.
